\documentclass[10pt,a4paper]{amsart}
\usepackage{amsmath,amssymb,amsthm,graphicx}
\usepackage{hyperref}
\usepackage{tikz}
%\usepackage{xcolor}
\usepackage{pict2e}
\let\arc\relax
%\usepackage{harpoon}
\usepackage{curves}
%\usepackage{chngcntr}
%\usepackage{apptools}
\usepackage{enumerate}

%\usepackage{verbatim}
%\usepackage[nospace,noadjust]{cite}
\usepackage{accents}
\usetikzlibrary{calc,decorations.markings}
\unitlength1cm

%%%%%%%%%%%%%%FONTS%%%%%%%%%%%%%%%%%%%%%%%%%%%%%%
\newcommand{\R}{\mathbb{R}}
\newcommand{\Z}{\mathbb{Z}}
\newcommand{\N}{\mathbb{N}}
\newcommand{\C}{\mathbb{C}}
\newcommand{\Id}{I}

%%%%%%%%%%%%%%MATHCAL%%%%%%%%%%%%%%%%%%%%%%%
\newcommand{\cK}{\mathcal{K}}
\newcommand{\cE}{\mathcal{E}}
\newcommand{\cR}{\mathcal{R}}
\newcommand{\cW}{\mathcal{W}}

%%%%%%%%%%%%%%%%%%ABBRS%%%%%%%%%%%%%%%%%%%%%%%%%%%%%

%%%%%%% OBJECTS SPECIFIC TO THIS PAPER%%%%%%%%%

\newcommand{\hw}{\widehat w}
\newcommand{\tw}{\widetilde w}

\newcommand{\bQ}{ Q}
\newcommand{\hQ}{{\widehat Q}}
\newcommand{\tQ}{{\widetilde Q}}

\newcommand{\sbQ}{Q^\star}
\newcommand{\shQ}{{\widehat Q^\star}}
\newcommand{\stQ}{{\widetilde Q^\star}}

\newcommand{\smu}{\mu^\star}
\newcommand{\shmu}{\widehat \mu^\star}
\newcommand{\stmu}{\widetilde \mu^\star}

\newcommand{\bmu}{{\mu}}
\newcommand{\hmu}{{\widehat \mu}}
\newcommand{\tmu}{{\widetilde \mu}}


\newcommand{\bP}{P}
\newcommand{\hP}{{\widehat P}}
\newcommand{\tP}{{\widetilde P}}

\newcommand{\sbP}{P^\star}
\newcommand{\shP}{{\widehat P^\star}}
\newcommand{\stP}{{\widetilde P^\star}}

\newcommand{\hR}{{\widehat R}}

\newcommand{\bv}{v}
\newcommand{\hv}{{\widehat v}}
\newcommand{\tv}{{\widetilde v}}

\newcommand{\sbv}{v^\star}
\newcommand{\shv}{{\widehat v^\star}}
\newcommand{\stv}{{\widetilde v^\star}}

\newcommand{\stu}{{\widetilde u^\star}}

\newcommand{\hF}{\widehat F}
\newcommand{\hW}{\widehat W}

\newcommand{\Fw}{\frac{F^2}{w}}
\newcommand{\hFw}{\frac{\hF^2}{\hw}}

\newcommand{\Cau}{\mathcal C_{\Sigma}}
\newcommand{\Cauv}{\mathcal C_{v}}
\newcommand{\Cauhv}{\mathcal C_{\hv}}
\newcommand{\Cautv}{\mathcal C_{\tv}}

\newcommand{\sCau}{\mathcal C_{\Sigma^\star}}
\newcommand{\sCauv}{\mathcal C_{\bv^\star}}
\newcommand{\sCauhv}{\mathcal C_{\hv^\star}}
\newcommand{\sCautv}{\mathcal C_{\stv}}

\newcommand{\Ltwo}{L^2(\Sigma)}
\newcommand{\sLtwo}{L^2(\Sigma^\star)}

\newcommand{\Sigl}{\Sigma^\ell}
\newcommand{\Sigr}{\Sigma^r}

%%%%%%%%%%%%%%%%OTHER SHORTCUTS%%%%%%%%%%%
\newcommand{\nn}{\nonumber}
\newcommand{\ul}{\underline}
\newcommand{\ol}{\overline}
\newcommand{\pa}{\partial}
\newcommand{\ti}{\tilde}
\newcommand{\wti}{\widetilde}
\newcommand{\wha}{\widehat}
\newcommand{\norm}[1]{\lVert#1 \rVert}
\newcommand{\abs}[1]{\lvert#1 \rvert}
%\newcommand{\Id}{\mathbb{I}}
\newcommand{\I}{\mathrm{i}}
\newcommand{\E}{\mathrm{e}}
\newcommand{\re}{\mathop{\mathrm{Re}}}
\newcommand{\im}{\mathop{\mathrm{Im}}}
\newcommand{\eps}{\varepsilon}
\newcommand{\dist}{\mathop{\mathrm{dist}}}

\newtheorem{Theorem}{Theorem}[section]
\newtheorem{Corollary}[Theorem]{Corollary}
\newtheorem{Lemma}[Theorem]{Lemma}
\newtheorem{Proposition}[Theorem]{Proposition}
 { \theoremstyle{definition}
\newtheorem{Definition}[Theorem]{Definition}
\newtheorem{Note}[Theorem]{Note}
\newtheorem{Example}[Theorem]{Example}
\newtheorem{Remark}[Theorem]{Remark} }

\numberwithin{equation}{section}
\usepackage{cleveref}
\crefname{originalreference}{original reference}{original references}
\Crefname{originalreference}{Original reference}{Original references}
\makeatletter
\newlabel{ContourGamma}{{internal}{source}{}{sourcepointer-1}{}}
\newlabel{ContourGamma@cref}{{[originalreference][0][]internal}{[1][0][]source}{}{}{}}
\makeatother
\title{Asymptotics of the squared Szegoe function divided by the weight at the two endpoints}
\author{Deift, Percy, Piorkowski, Mateusz}
\date{doi:10.3842/SIGMA.2024.004; CC BY 4.0}
\begin{document}
\maketitle
\setcounter{section}{2}
\setcounter{Theorem}{1}
\noindent\textit{Editorial note.} All delivered mathematical text is the author's own, copied from the pinned SIGMA article source; no statement or proof step was rewritten, reordered or removed. The counted statement is the labeled proposition environment at source lines 438-452 and the counted proof is the proof environment at source lines 474-510.\par\medskip
\noindent\textit{Editorial note.} The proof attributes the first asymptotic formula of the statement to a published proposition of the cited literature and derives the second one; the attribution is preserved verbatim and the cited reference is declared as a standard prerequisite of this card.\par\medskip
\noindent\textit{Editorial note.} The proof cites published references of the author's bibliography by key; the cited entries are copied line for line from the author's own bibliography into a retained bibliography fragment, and hyperref is added because the publisher class, not the author preamble, supplied the link command.\par\medskip
\noindent\textit{Editorial note.} The retained preamble of this unit was generated mechanically for this delivery: it copies the author's own package, numbering, macro and theorem declaration block from the pinned source (source lines 9-133), with the publisher class and page setup lines removed.\par\medskip
\noindent\textit{Allowed standard prerequisites (editorial).}\begin{itemize}
\item The asymptotic expansion of the squared Szegoe function divided by the weight at the endpoint one. published statement cited by the author at proof line 475. Supplies the first asymptotic formula of the statement, which the author quotes instead of rederiving it.
\item The explicit integral representation of the Szegoe function for a weight on the interval. derived inside the proof at proof lines 477-486. Provides the representation whose exponential argument is analysed to obtain the behaviour at the endpoint minus one.
\item The residue computation transferring a Cauchy integral along the cut to a contour integral. derived inside the proof at proof lines 489-497. Rewrites the exponential argument in the form that exhibits the constant defined by the contour integral in the statement.
\end{itemize}\medskip
\ifdefined\hypertarget\hypertarget{sourcepointer-1}{}\fi
\noindent\textit{Source reference.} Original reference internal: sources/sigma-2024-004-recurrence-coefficients-for-orthogonal-polyn/pinned-original-source.tex, source lines 471--471. The proof refers to the contour depicted in the figure between the statement and the proof by its caption label; that figure is the author's own and is not part of either counted span, and every step of the counted proof is retained verbatim.\par\medskip
\begin{Proposition} \label{F^2/w}
The function $\frac{F^2}{w} \colon \C \setminus [-1,\infty) \to \C$ satisfies
\begin{align}\label{F^2/w at +1}
 \frac{F^2}{w}(z) = 1 \mp \frac{\I \pi}{w(z)}-\frac{\pi^2}{2w^2(z)} +O\left(\frac{1}{\log^3(z-1)}\right)
\end{align}
uniformly for any path $z\to +1$ in $\C \setminus [-1,\infty)$, where the $-$ is taken for~$\im(z)> 0$ and $+$ for~$\im(z) < 0$, and also
\begin{align}\label{F^2/w at -1}
 \frac{F^2}{w}(z) = \phi(z)^{-1}\exp\bigl\lbrace-\big(z^2-1\big)^{1/2}\big(d_0+O(|z+1|)\big)\bigr\rbrace,
\end{align}
uniformly for any path $z\to-1$ in $\C \setminus [-1,\infty)$, where
\begin{align}\label{d0}
 d_0 = \frac{1}{2\pi\I} \int_\gamma \frac{\log\big(w(\zeta)/(1+\zeta)\big)}{(\zeta^2-1)^{1/2}}\frac{{\rm d}\zeta}{\zeta+1}
\end{align}
and $\gamma$ is an oriented contour originating from the point $\zeta = 1$, going anticlockwise around the interval $[-1,1]$ and ending again at the point $\zeta = 1$ as depicted in Figure~{\rm \ref{ContourGamma}}.
\end{Proposition}

\begin{proof}
For the proof of statement \eqref{F^2/w at +1}, see \cite[Proposition~A.1]{CD}. Statement \eqref{F^2/w at -1} is a special case of \cite[Lemma~6.6]{KMVV}, but with an additional restriction on the choice of the contour $\gamma$ stemming from the logarithmic singularity of $w$. Due to this technicality, we repeat the proof found therein.

First, let us note that if $F_1$ is the Szeg\H{o} function of a weight $w_1$ and $F_2$ the Szeg\H{o} function of a weight $w_2$, then the Szeg\H{o} function $F_{12}$ of the product $w_1w_2$ is given by the product of the individual Szeg\H{o} functions: $F_{12} = F_1 F_2$. As the Szeg\H{o} function of a Jacobi weight $w_{\alpha,\beta}(x) = (1-x)^\alpha (1+x)^\beta$ with $\alpha, \beta > -1$, is given by (see \cite[Remark 5.1]{KMVV})
\begin{align*}
 F_{\alpha,\beta}(z) = \frac{(z-1)^{\alpha/2}(z+1)^{\beta/2}}{\phi(z)^{\frac{\alpha+\beta}{2}}},
\end{align*}
we can conclude that
\begin{align}
 F(z) ={}& \frac{(z+1)^{1/2}}{\phi(z)^{1/2}}\nonumber\\
 &\times\exp \Bigg( \dfrac{\big(z^2-1\big)^{1/2}}{2 \pi}\int_{-1}^1 \dfrac{\log (w(s)/(1+s))}{\sqrt{1-s^2}} \dfrac{{\rm d}s}{z-s} \Bigg),
 \qquad z \in \C \setminus [-1,1].\label{SzegoFact}
\end{align}
Again, $(z+1)^{1/2}$ and $\phi(z)^{1/2}$ denote the principal branches (a simple calculation shows that $\phi(z) \not \in (-\infty, -1]$ for $z \in \C \setminus (-\infty, 1]$). Moreover, for $x < -1$, $\big(\phi(x)^{1/2}\big)_+ = -\big(\phi(x)^{1/2}\big)_-$ as~$\phi(z) \sim 2z$ for $z \to \infty$. Thus $(z+1)^{1/2}/\phi(z)^{1/2}$ is indeed analytic in $\C\setminus [-1,1]$.

To analyze the argument of the exponential in \eqref{SzegoFact}, we choose for a fixed $z \in \C \setminus [-1,1]$ a~contour $\gamma$ as shown in Figure~\ref{ContourGamma}. Then a residue calculation shows that
\begin{align}\nonumber
 &\frac{1}{\pi} \int_{-1}^1 \frac{\log(w(s)/(1+s))}{\sqrt{1-s^2}} \frac{{\rm d}s}{z-s}
 \\\label{ResiCalc}
 &\qquad= \frac{\log(w(z)/(1+z))}{\big(z^2-1\big)^{1/2}} - \frac{1}{2\pi\I} \int_\gamma\frac{\log(w(\zeta)/(1+\zeta))}{\big(\zeta^2-1\big)^{1/2}} \frac{{\rm d}\zeta}{\zeta-z}.
\end{align}
Here we note the $w(z)/(1+z)$ is analytic and non-zero in the simply connected region $\C \setminus [1,\infty)$, and hence $\log(w(z)/(1+z))$ exists and is analytic in this region. Moreover, the integrand in~\eqref{ResiCalc} is integrable along $\gamma$ as long as $z \not \in \gamma$.

Note also that as $\log(w(\zeta)/(1+\zeta))$ has an iterated logarithmic singularity at $\zeta = 1$, we cannot deform $\gamma$ away from this point. Plugging \eqref{SzegoFact} and \eqref{ResiCalc} into the definition of $\Fw$, we obtain
\begin{align*}
 \Fw(z) = \phi(z)^{-1} \exp \Biggl( - \frac{\big(z^2-1\big)^{1/2}}{2\pi\I} \int_\gamma\frac{\log(w(\zeta)/(1+\zeta))}{\big(\zeta^2-1\big)^{1/2}} \frac{{\rm d}\zeta}{\zeta-z} \Biggr).
\end{align*}
We compute the Taylor expansion around $\zeta = -1$ to obtain
\begin{align*}
 \frac{1}{2\pi\I} \int_\gamma\frac{\log(w(\zeta)/(1+\zeta))}{\big(\zeta^2-1\big)^{1/2}} \frac{{\rm d}\zeta}{\zeta-z} = \sum_{k=0}^\infty d_k (z+1)^k,
\end{align*}
where $d_k$ is given by
\begin{align*}
 d_k = \frac{1}{2\pi\I} \int_\gamma \frac{\log\big(w(\zeta)/(1+\zeta)\big)}{(\zeta^2-1)^{1/2}}\frac{{\rm d}\zeta}{(\zeta+1)^{k+1}}.
\end{align*}
This finishes the proof.
\end{proof}

\begin{thebibliography}{99}
\bibitem{CD}
Conway T.O., Deift P., Asymptotics of polynomials orthogonal with respect to a
 logarithmic weight, \href{https://doi.org/10.3842/SIGMA.2018.056}{\textit{SIGMA}} \textbf{14} (2018), 056, 66~pages,
 \href{https://arxiv.org/abs/1711.01590}{arXiv:1711.01590}.


\bibitem{KMVV}
Kuijlaars A.B.J., McLaughlin K.T.-R., Van~Assche W., Vanlessen M., The
 {R}iemann--{H}ilbert approach to strong asymptotics for orthogonal
 polynomials on {$[-1,1]$}, \href{https://doi.org/10.1016/j.aim.2003.08.015}{\textit{Adv. Math.}} \textbf{188} (2004), 337--398,
 \href{https://arxiv.org/abs/math.CA/0111252}{arXiv:math.CA/0111252}.
\end{thebibliography}
\end{document}
